% source: https://github.com/pfradin/luadraw/discussions/110#discussioncomment-14559937 (pfradin, block 1)
% !TEX TS-program = lualatex

%%%
% src::
%     url = https://tex.stackexchange.com/a/751717/6880
%%%

\documentclass{standalone}

\usepackage[svgnames]{xcolor}
\usepackage[3d]{luadraw}

\begin{document}

\begin{luadraw}{name = spherical-triangle}
require 'luadraw_spherical'

local graphview = graph3d:new{
  window  = {-4, 4, -3, 5},
  viewdir = {25, 70},
  size    = {10, 10}
}

graphview:Linewidth(6)
Hiddenlinestyle = "dashed"
Hiddenlines     = true

arrowBstyle = "-stealth"

local O, R =  Origin, 3
-- Caution: you must always define the sphere first, the sM function for example, uses the sphere data
graphview:Define_sphere({center = O, radius = R})

local veci, vecj, veck = graphview:Proj3d(vecI), graphview:Proj3d(vecJ), graphview:Proj3d(vecK)
local kx, ky, kz = (R+1)/cpx.abs(veci), (R+1)/cpx.abs(vecj), (R+1)/cpx.abs(veck)
-- veci is a direction vector of the projection of the Ox axis and 
-- the projection of the 3d point O+kx*vecI is: o+kx*veci (o is the projection of O) 
-- and its distance to o is R+1

local A, B = sM(-50, 60), sM(-10, 60)
local C, D, N = sM(50, 70), sM(170, 70), sM(90, 0)

graphview:DSfacet(
  {N, A, B},
  {fill = "brown", fillopacity = 0.6}
)

graphview:DSfacet(
  {N, C, D},
  {fill = "brown", fillopacity = 0.6}
)

graphview:DSpolyline(
  {
    {O, O + kx * vecI},
    {O, O + ky * vecJ},
    {O, O + kz * vecK}},
    {arrows = 1, width = 6, color = "blue"}
)

graphview:DScircle(
  {O, vecK},       -- Plan du grand cercle.
  {color = "red"}
)

graphview:Dspherical()

graphview:Dlabel3d(
  "$N$", N, {pos = "NW"},
  "$O$", O, {pos = "S"},
  "$x$", O + kx * vecI, {pos = "SW"},
  "$y$", O + ky * vecJ , {pos = "E"},
  "$z$", O + kz * vecK , {pos = "N"},
  "$A$", A, {pos = "S"},
  "$B$", B, {},
  "$C$", C, {pos = "E"},
  "$D$", D, {pos = "S"}
)
graphview:Show()
\end{luadraw}
\end{document}
